\documentclass[12pt,a4paper]{article}

% ============================================================
%  PACKAGES
% ============================================================
\usepackage{fontspec}
\usepackage{polyglossia}
\setmainlanguage{vietnamese}
\setotherlanguage{english}
\setmainfont{FreeSerif}[
  BoldFont = FreeSerifBold,
  ItalicFont = FreeSerifItalic,
  BoldItalicFont = FreeSerifBoldItalic
]
\setsansfont{FreeSans}
\setmonofont{FreeMono}
\usepackage{amsmath,amsfonts,amssymb}
\usepackage{geometry}
\usepackage{hyperref}
\usepackage{booktabs}
\usepackage{array}
\usepackage{tabularx}
\usepackage{longtable}
\usepackage{multirow}
\usepackage[table]{xcolor}
\usepackage{enumitem}
\usepackage{setspace}
\usepackage{fancyhdr}
\usepackage{titlesec}
\usepackage{graphicx}
\usepackage{algorithm}
\usepackage{algpseudocode}
\usepackage{caption}
\usepackage{subcaption}
\usepackage{natbib}
\usepackage{tikz}
\usetikzlibrary{shapes.geometric, arrows.meta, positioning, fit}
\usepackage{pgfplots}
\pgfplotsset{compat=1.18}

% ============================================================
%  PAGE SETUP
% ============================================================
\geometry{
  a4paper,
  top=2.5cm, bottom=2.5cm,
  left=3.0cm, right=2.5cm
}
\setstretch{1.3}

% ============================================================
%  COLORS
% ============================================================
\definecolor{primary}{RGB}{31, 56, 100}
\definecolor{accent}{RGB}{46, 80, 144}
\definecolor{lightblue}{RGB}{235, 242, 252}
\definecolor{lightgreen}{RGB}{235, 252, 238}
\definecolor{lightyellow}{RGB}{255, 252, 235}
\definecolor{lightred}{RGB}{255, 240, 240}
\definecolor{darkgray}{RGB}{80, 80, 80}
\definecolor{codebg}{RGB}{248, 248, 248}

% ============================================================
%  SECTION FORMATTING
% ============================================================
\titleformat{\section}
  {\color{primary}\large\bfseries}
  {\thesection.}{0.6em}{}
  [\color{accent}\rule{\linewidth}{0.6pt}\vspace{2pt}]

\titleformat{\subsection}
  {\color{accent}\normalsize\bfseries}
  {\thesubsection.}{0.5em}{}

\titleformat{\subsubsection}
  {\color{darkgray}\normalsize\bfseries\itshape}
  {\thesubsubsection.}{0.5em}{}

% ============================================================
%  HEADER / FOOTER
% ============================================================
\pagestyle{fancy}
\fancyhf{}
\fancyhead[L]{\small\color{darkgray} Đề cương NCKH -- IRP-TW-DT Hà Nội}
\fancyhead[R]{\small\color{darkgray} \thepage}
\renewcommand{\headrulewidth}{0.4pt}

% ============================================================
%  CUSTOM COMMANDS
% ============================================================
\newcommand{\highlight}[1]{\colorbox{lightblue}{\parbox{\dimexpr\linewidth-2\fboxsep}{#1}}}
\newcommand{\boldterm}[1]{\textbf{\color{primary}#1}}
\newcommand{\citealias}[1]{\citeauthor{#1} (\citeyear{#1})}

% Math shortcuts
\newcommand{\T}{\mathcal{T}}
\newcommand{\N}{\mathcal{N}}
\newcommand{\K}{\mathcal{K}}
\newcommand{\V}{\mathcal{V}}
\newcommand{\A}{\mathcal{A}}
\newcommand{\vij}[1]{v_{ij}^{(#1)}}
\newcommand{\tij}{\tau_{ij}}

% ============================================================
%  DOCUMENT
% ============================================================
\begin{document}

% ============================================================
%  TITLE PAGE
% ============================================================
\begin{titlepage}
\begin{center}
  \vspace*{1.0cm}

  {\small\scshape Trường Đại học Bách Khoa Hà Nội}\\[0.3cm]
  {\small\scshape Viện Toán ứng dụng và Tin học -- Khoa Khoa học Máy tính}\\[1.0cm]

  \rule{\linewidth}{1.5pt}\\[0.4cm]

  {\LARGE\bfseries\color{primary}
    ĐỀ CƯƠNG NGHIÊN CỨU KHOA HỌC}\\[0.6cm]

  {\large\bfseries\color{accent}
    TỐI ƯU HÓA HỆ THỐNG GIAO HÀNG CHẶNG CUỐI TRONG \\
    LOGISTICS ĐÔ THỊ BẰNG MÔ HÌNH INVENTORY ROUTING\\
    PROBLEM KẾT HỢP KHUNG THỜI GIAN VÀ\\
    GIAO THÔNG ĐỘNG}\\[0.4cm]

  {\normalsize\itshape\color{darkgray}
    (IRP with Time Windows and Dynamic Traffic -- IRP-TW-DT)}\\[0.2cm]

  \rule{\linewidth}{1.5pt}\\[1.5cm]

  \begin{flushleft}
    \begin{tabular}{@{}ll}
      \textbf{Sinh viên thực hiện:}  &  \\[4pt]
      \textbf{Mã sinh viên:}          & [Mã SV] \\[4pt]
      \textbf{Giảng viên hướng dẫn:} & [Tên Giảng viên] \\[4pt]
      \textbf{Chuyên ngành:}          & Khoa học Máy tính / \\
                                      & Operations Research \& Logistics Optimization \\[4pt]
      \textbf{Năm học:}               & 2025--2026 \\
    \end{tabular}
  \end{flushleft}

  \vfill
  {\normalsize\itshape Hà Nội, Tháng 3 Năm 2026}
\end{center}
\end{titlepage}

% ============================================================
%  TABLE OF CONTENTS
% ============================================================
\newpage
\tableofcontents
\newpage

% ============================================================
\section{Giới thiệu và Tính cấp thiết}
% ============================================================

Trong chuỗi cung ứng hiện đại, chi phí logistics chặng cuối (\textit{last-mile delivery})
thường chiếm tỷ trọng lên tới 40--50\% tổng chi phí vận chuyển \citep{gevaers2011characterization},
và có thể vượt quá 50\% trong các thị trường châu Á--Thái Bình Dương có mật độ đô thị cao
\citep{fedex2024apac}. Tại bối cảnh Việt Nam, chi phí logistics chiếm khoảng 25\% GDP ---
cao hơn đáng kể so với mức 15\% ở các nước phát triển \citep{amr2022vietnam}.
Tại Hà Nội --- đô thị hơn 8 triệu dân với đặc trưng giao thông hỗn hợp và mật độ xe
máy cực cao --- thách thức vận tải càng được khuếch đại. Theo TomTom Traffic Index 2025
\citep{tomtom2025hanoi}, mức độ tắc nghẽn trung bình cả ngày tại Hà Nội đạt 49,2\%
(tức hành trình thực tế dài hơn điều kiện lý tưởng gần 50\%); riêng giờ cao điểm sáng
tốc độ trung bình chỉ đạt 15,3~km/h và giờ cao điểm chiều còn tệ hơn với 13,6~km/h.

Hầu hết các doanh nghiệp logistics Việt Nam hiện đang vận hành theo mô hình
\boldterm{Retailer Managed Inventory (RMI)}: khách hàng tự đặt hàng khi tồn kho cạn,
và đội xe phản ứng theo yêu cầu. Cách tiếp cận này dẫn đến hai hệ quả bất lợi đồng thời:
(1)~chuyến giao thường bị buộc vào những khung giờ kém linh hoạt, thường trùng với giờ
cao điểm; (2)~không tận dụng được việc tổng hợp nhiều đơn hàng nhỏ thành một chuyến
tối ưu.

Bài toán \boldterm{Inventory Routing Problem (IRP)} được đề xuất như giải pháp căn bản:
thay vì để khách hàng tự quản lý tồn kho, nhà cung cấp chủ động quyết định \textit{giao
cho ai, giao bao nhiêu và đi theo lộ trình nào} trên chân trời kế hoạch nhiều ngày.
IRP là sự tích hợp giữa quản lý tồn kho và bài toán định tuyến phương tiện có ràng buộc
(CVRP), thuộc lớp NP-Hard \citep{coelho2014thirty}.

Tuy nhiên, tuyệt đại đa số các mô hình IRP trong văn học hiện tại giả định thời gian
di chuyển là hằng số. \citealias{touzout2022assign} là nhóm đầu tiên đề xuất mô hình
IRP với thời gian di chuyển phụ thuộc thời điểm xuất phát (TD-IRP), nhưng công trình
này chưa tích hợp ràng buộc khung thời gian phục vụ (\textit{time windows}) và chưa
được kiểm chứng trong bối cảnh đô thị châu Á với đặc trưng giao thông riêng biệt.

Nghiên cứu này đề xuất mô hình \boldterm{IRP-TW-DT} (IRP with Time Windows and Dynamic
Traffic) --- mở rộng của TD-IRP với bốn đóng góp cụ thể:

\begin{enumerate}[leftmargin=*, label=(\arabic*)]
  \item Mô hình hóa giao thông động Hà Nội bằng \textit{piecewise constant speed function}
        đảm bảo tính chất FIFO, calibrate theo múi giờ thực tế của thành phố;
  \item Tích hợp \textit{time windows} $[e_{it}, l_{it}]$ cho từng khách hàng trong
        từng ngày, phản ánh thực tế vận hành của các cửa hàng và đại lý đô thị;
  \item Thiết kế thuật toán \textit{Hybrid Genetic Algorithm} (HGA) với cơ chế mã hóa
        hai tầng (Two-part Chromosome) và toán tử \textit{Time-Shift} mới;
  \item Đánh giá định lượng sự đánh đổi giữa chi phí lưu kho và chi phí vận tải dưới
        tác động của giao thông động trong môi trường đô thị.
\end{enumerate}

% ============================================================
\section{Tổng quan tài liệu}
% ============================================================

\subsection{Bài toán Inventory Routing Problem (IRP)}

Bài toán IRP có nguồn gốc từ mô hình phân phối khí công nghiệp của Bell et al. (1983)
và được tổng hợp toàn diện bởi \citealias{coelho2014thirty} trong 30 năm phát triển.
IRP tìm kiếm chính sách phân phối tối ưu theo mô hình
\boldterm{Vendor Managed Inventory (VMI)}: nhà cung cấp quản lý tồn kho tại khách hàng
và tự quyết định thời điểm, số lượng giao hàng. \citealias{archetti2016vmi} chứng minh
rằng VMI có thể tiết kiệm trung bình 24\% so với RMI, song kết quả này dựa trên giả
định thời gian di chuyển hằng số.

Các phương pháp giải chính xác (Branch-and-Cut, Branch-and-Price) như của
\citealias{archetti2007branch} và \citealias{desaulniers2016branch} chỉ giải quyết
hiệu quả các instance dưới 30 khách hàng. Đối với quy mô lớn hơn, các Metaheuristic
như Large Neighborhood Search (LNS) hay HGA được ưu tiên
\citep{bouvier2024continent, adulyasak2014branch}.

\subsection{Bài toán Định tuyến Phụ thuộc Thời gian (TDVRP)}

\citealias{ichoua2003vehicle} đề xuất mô hình tốc độ hằng từng đoạn (\textit{piecewise
constant speed function}) cho TDVRP. Mô hình IGP chia ngày thành $S$ khoảng thời gian
$[a_s, a_{s+1})$, mỗi khoảng gán một tốc độ hằng $v^{(s)}$ cho mọi cung trên đồ thị.
\citealias{ghiani2014note} chứng minh rằng mô hình này là trường hợp tổng quát của
mọi hàm di chuyển tuyến tính từng đoạn thỏa mãn tính chất FIFO và hiện là tiêu chuẩn
de-facto trong toàn bộ văn học TDVRP hiện đại \citep{adamo2024review}.

Trong bối cảnh đô thị, \citealias{rifki2020impact} chứng minh sự granularity theo cả
không gian lẫn thời gian của dữ liệu traffic ảnh hưởng đáng kể đến chất lượng lời giải,
nhấn mạnh tầm quan trọng của việc calibrate mô hình từ dữ liệu thực địa.

\subsection{Bài toán IRP Phụ thuộc Thời gian (TD-IRP)}

Đây là hướng nghiên cứu mới nhất và trực tiếp liên quan đến đề tài. Bảng~\ref{tab:related}
tổng hợp các công trình nổi bật.

\begin{table}[htbp]
\centering
\caption{Các công trình IRP liên quan đến giao thông động}
\label{tab:related}
\renewcommand{\arraystretch}{1.35}
\begin{tabularx}{\linewidth}{@{}l l l c c c c@{}}
\toprule
\textbf{Tác giả} & \textbf{Tạp chí} & \textbf{Năm} &
\textbf{TD} & \textbf{TW} & \textbf{Đa xe} & \textbf{Metah.} \\
\midrule
Archetti et al.       & Transp. Sci.    & 2007 & -- & -- & \checkmark & -- \\
Coelho et al.         & Transp. Sci.    & 2014 & -- & -- & \checkmark & \checkmark \\
Touzout et al.        & EJOR            & 2022 & \checkmark & -- & \checkmark & \checkmark \\
Skålnes et al.        & EJOR            & 2022 & -- & -- & \checkmark & -- \\
Bouvier et al.        & Transp. Sci.    & 2024 & -- & -- & \checkmark & \checkmark \\
Skålnes et al.        & EJOR            & 2024 & -- & -- & \checkmark & \checkmark \\
\midrule
\rowcolor{lightblue}
\textbf{Đề tài này}   & --              & 2026 & \checkmark & \checkmark & \checkmark & \checkmark \\
\bottomrule
\end{tabularx}
\vspace{3pt}
\small\textit{TD: Time-Dependent travel time; TW: Time Windows; Metah.: Metaheuristic}
\end{table}

\boldterm{Touzout et al. (2022) \citep{touzout2022assign}} --- công trình quan trọng nhất
cần làm rõ vị thế của đề tài --- đề xuất mô hình TD-IRP và thuật toán Matheuristic hai
pha (phân công khách hàng trước, định tuyến sau) được thử nghiệm trên bộ dữ liệu thực
từ thành phố Lyon, Pháp. Tuy nhiên, công trình này có ba giới hạn chính:

\begin{enumerate}[leftmargin=*, label=(\arabic*)]
  \item \textit{Không có time windows}: Touzout et al. không ràng buộc khung giờ nhận hàng
        cho khách --- một thực tế phổ biến trong logistics đô thị.
  \item \textit{Bối cảnh đô thị châu Âu}: Profile giao thông Lyon, Pháp khác biệt căn bản
        so với Hà Nội (mạng lưới đường, mật độ xe máy, đặc trưng giờ cao điểm kép).
  \item \textit{Thuật toán Matheuristic (Tabu + MIP)}: Đề tài này đóng góp một thuật toán
        hoàn toàn khác là HGA với Two-part Chromosome và toán tử Time-Shift mới.
\end{enumerate}

\subsection{Khoảng trống Nghiên cứu}

Qua tổng quan tài liệu, tồn tại một khoảng trống nghiên cứu được xác định rõ:

\begin{center}
\setlength{\fboxsep}{10pt}
\fbox{\begin{minipage}{0.92\linewidth}
\textbf{Chưa tồn tại mô hình IRP-TW-DT:} Mặc dù \citealias{touzout2022assign} đã xây
dựng nền tảng TD-IRP, sự kết hợp của \emph{(i)}~time windows, \emph{(ii)}~profile giao
thông đô thị nhiều đỉnh calibrate từ dữ liệu thực, và \emph{(iii)}~thuật toán HGA với
cơ chế tối ưu theo chiều thời gian chưa từng được nghiên cứu. Đặc biệt, chưa có công
trình nào định lượng sự đánh đổi giữa chi phí lưu kho và chi phí vận tải dưới tác động
của giao thông động khi chuyển đổi từ RMI sang VMI.
\end{minipage}}
\end{center}

% ============================================================
\section{Câu hỏi nghiên cứu}
% ============================================================

\begin{enumerate}[leftmargin=*, label=\textbf{Q\arabic{*}.}, itemsep=6pt]
  \item Giao thông đô thị phụ thuộc thời gian ảnh hưởng như thế nào đến sự đánh đổi
        giữa chi phí tồn kho và chi phí vận tải trong bài toán IRP?

  \item Thuật toán HGA với cấu trúc Two-part Chromosome (Allocation + Routing)
        có thể đồng thời tối ưu hóa quyết định tồn kho và lịch định tuyến
        trong khung IRP-TW-DT không?

  \item Time-Shift neighborhood search cải thiện chất lượng lời giải như thế nào
        so với HGA không có local search trên trục thời gian?

  \item Mô hình IRP-TW-DT tiết kiệm được bao nhiêu phần trăm tổng chi phí so với
        các baseline (Periodic Delivery, RMI) trên bộ dữ liệu mô phỏng Hà Nội?
\end{enumerate}

% ============================================================
\section{Mục tiêu nghiên cứu}
% ============================================================

\begin{enumerate}[leftmargin=*, label=\textbf{M\arabic*}.\ , itemsep=8pt]
  \item \boldterm{Mục tiêu mô hình hóa:}
    Xây dựng mô hình MILP cho bài toán IRP-TW-DT đa chu kỳ, trong đó thời gian di
    chuyển phụ thuộc thời điểm xuất phát theo profile giao thông Hà Nội 5 múi giờ.
    Đảm bảo tính chất FIFO theo chuẩn \citep{ichoua2003vehicle}.

  \item \boldterm{Mục tiêu thuật toán:}
    Thiết kế HGA với cơ chế mã hóa Two-part Chromosome và toán tử Time-Shift để giải
    IRP-TW-DT trong thời gian khả thi; đồng thời validate HGA trên bộ benchmark chuẩn
    của \citealias{archetti2007branch}.

  \item \boldterm{Mục tiêu thực nghiệm:}
    Đánh giá và so sánh bốn kịch bản (Periodic / RMI / IRP-TW-Static / IRP-TW-DT)
    trên bộ dữ liệu mô phỏng mạng lưới phân phối Hà Nội (quy mô 20, 50, 100 khách
    hàng, 5 seeds) nhằm rút ra hàm ý quản trị cho doanh nghiệp logistics Việt Nam.
\end{enumerate}

% ============================================================
\section{Mô hình Toán học của bài toán IRP-TW-DT}
% ============================================================

\subsection{Mô hình Giao thông Động (Traffic Model)}

\subsubsection{Piecewise Constant Speed Function (IGP Model)}

Theo chuẩn \citep{ichoua2003vehicle, ghiani2014note}, ngày được chia thành $S$ khoảng
thời gian (\textit{time zones}) $[a_1, a_2), [a_2, a_3), \ldots, [a_S, a_{S+1})$ với
$a_1 = 0$ (0h) và $a_{S+1} = H$ (24h). Tốc độ di chuyển trên mọi cung là hằng số
trong mỗi khoảng:

\begin{equation}
  v^{(s)} = \text{tốc độ } [\text{km/h}] \text{ trong khoảng } [a_s, a_{s+1}),
  \quad s = 1, \ldots, S
\end{equation}

\noindent Bảng~\ref{tab:speed} trình bày profile 5 múi giờ calibrate theo đặc trưng
Hà Nội (cần validate bằng dữ liệu thực, ví dụ: Google Maps Distance Matrix API hoặc
báo cáo của TEDI/TomTom).

\begin{table}[htbp]
\centering
\caption{Profile giao thông Hà Nội --- Piecewise Constant Speed (hiệu chỉnh từ TomTom 2025)}
\label{tab:speed}
\renewcommand{\arraystretch}{1.2}
\begin{tabular}{@{}c c c c@{}}
\toprule
\textbf{Zone} & \textbf{Múi giờ} & \textbf{Tốc độ $v^{(s)}$ (km/h)} &
\textbf{Nhận xét} \\
\midrule
1 & $[0\text{h},\ 6\text{h})$    & 27 & Đêm, giao thông thưa \\
2 & $[6\text{h},\ 9\text{h})$    & 15 & Cao điểm sáng \\
3 & $[9\text{h},\ 17\text{h})$   & 19 & Hành chính \\
4 & $[17\text{h},\ 19\text{h})$  & 14 & Cao điểm chiều (nặng nhất) \\
5 & $[19\text{h},\ 24\text{h})$  & 21 & Tối, giao thông dịu \\
\bottomrule
\end{tabular}
\vspace{4pt}
\begin{minipage}{0.9\linewidth}
\small\textit{Nguồn: Hiệu chỉnh từ TomTom Traffic Index 2025 \citep{tomtom2025hanoi} ---
dữ liệu thực tế: tốc độ giờ cao điểm sáng 15,3~km/h, chiều 13,6~km/h; thời gian di chuyển
trung bình 10~km là 33~phút 20~giây. Các giá trị trong bảng là ước tính tốt nhất từ
dữ liệu công khai; cần hiệu chỉnh lại bằng dữ liệu GPS thực địa hoặc Google Maps
Distance Matrix API trước khi sử dụng trong thực nghiệm chính thức.}
\end{minipage}
\end{table}

\subsubsection{Tính toán Thời gian Di chuyển và Tính chất FIFO}

Thời gian di chuyển $\tij(\tau)$ từ điểm $i$ đến điểm $j$ khi xuất phát lúc $\tau$
được tính bằng thuật toán tích phân rời rạc qua các time zones:

\begin{equation}
  \tij(\tau) = \int_{\tau}^{\tau + \tij(\tau)} \frac{d_{ij}}{v(\xi)} \, d\xi
  \label{eq:travel_time}
\end{equation}

\noindent Trong thực tế, với speed function bậc thang, phương trình (\ref{eq:travel_time})
được giải bằng cách cộng dồn quãng đường qua từng zone:

\begin{equation}
  \tij(\tau) = \sum_{s=s_\tau}^{S} \min\!\left(\frac{r_s}{v^{(s)}},\,
                                    a_{s+1} - \max(\tau, a_s)\right)
  \label{eq:tt_discrete}
\end{equation}

\noindent trong đó $r_s$ là khoảng cách còn lại khi xe bước vào zone $s$, và $s_\tau$
là zone chứa $\tau$. Mô hình IGP đảm bảo tính chất FIFO: với mọi $\tau_1 < \tau_2$,
luôn có $\tau_1 + \tij(\tau_1) \leq \tau_2 + \tij(\tau_2)$
\citep{ghiani2014note}.

\subsection{Ký hiệu và Tập hợp}

\begin{description}[leftmargin=3cm, labelwidth=2.5cm, font=\normalfont\bfseries]
  \item[$\T = \{1,\ldots,T\}$] Tập các chu kỳ (ngày) trong chân trời kế hoạch.
  \item[$\N' = \{1,\ldots,n\}$] Tập khách hàng. Node $0$ là depot.
  \item[$\V = \N' \cup \{0\}$] Tập tất cả các điểm.
  \item[$\A$] Tập các cung $(i,j)$ với $i \neq j$, $i,j \in \V$.
  \item[$\K = \{1,\ldots,m\}$] Tập phương tiện đồng nhất.
\end{description}

\subsection{Tham số}

\begin{description}[leftmargin=4cm, labelwidth=3.5cm, font=\normalfont\itshape]
  \item[$h_i$] Chi phí lưu kho đơn vị tại khách hàng $i$ (đ/đv/ngày).
  \item[$U_i$] Dung tích kho tối đa tại $i$.
  \item[$L_i^{\min}$] Mức tồn kho an toàn tối thiểu tại $i$.
  \item[$I_{i0}$] Tồn kho ban đầu tại $i$.
  \item[$d_{it}$] Nhu cầu tiêu thụ tất định của khách $i$ trong ngày $t$.
  \item[$Q$] Tải trọng xe (đội xe đồng nhất).
  \item[$e_i,\ l_i$] Cận dưới và cận trên của khung thời gian phục vụ tại khách $i$
                      (giờ trong ngày; giống nhau mọi ngày --- giả định tiêu chuẩn).
  \item[$s_i$] Thời gian phục vụ (bốc xếp hàng) tại điểm $i$ (giờ); $s_i = 0.25$h
               (15 phút) cho mọi khách trong thực nghiệm.
  \item[$d_{ij}$] Khoảng cách Euclidean giữa điểm $i$ và $j$ (km).
  \item[$\tij(\tau)$] Thời gian di chuyển từ $i$ đến $j$ khi xuất phát lúc $\tau$ (giờ).
  \item[$c_d$] Chi phí đơn vị theo khoảng cách (đ/km).
  \item[$c_t$] Chi phí đơn vị theo thời gian (đ/giờ; lương lái xe + nhiên liệu theo giờ).
  \item[$M$] Hằng số Big-M đủ lớn.
\end{description}

\subsection{Biến quyết định}

\begin{description}[leftmargin=4cm, labelwidth=3.5cm, font=\normalfont\itshape]
  \item[$I_{it} \geq 0$] Mức tồn kho tại khách $i$ cuối ngày $t$.
  \item[$q_{ikt} \geq 0$] Lượng hàng giao cho khách $i$ bởi xe $k$ trong ngày $t$.
  \item[$x_{ijkt} \in \{0,1\}$] 1 nếu xe $k$ đi từ $i$ đến $j$ trong ngày $t$.
  \item[$z_{ikt} \in \{0,1\}$] 1 nếu xe $k$ ghé thăm khách $i$ trong ngày $t$.
  \item[$w_{ikt} \geq 0$] Thời điểm xe $k$ bắt đầu phục vụ tại $i$ ngày $t$
                           (với $i \in \N'$).
  \item[$w_{0kt} \geq 0$] Thời điểm xe $k$ xuất phát từ depot trong ngày $t$,
                           $w_{0kt} \leq H$ (biến tự do, không bị ràng buộc time window).
\end{description}

\subsection{Hàm mục tiêu}

Tối thiểu hóa tổng chi phí, bao gồm ba thành phần riêng biệt:

\begin{equation}
  \min\ Z = \underbrace{\sum_{t \in \T} \sum_{i \in \N'} h_i \cdot I_{it}}_{\text{Chi phí lưu kho}}
        + \underbrace{\sum_{t \in \T} \sum_{k \in \K} \sum_{(i,j) \in \A}
                c_d \cdot d_{ij} \cdot x_{ijkt}}_{\text{Chi phí khoảng cách}}
        + \underbrace{\sum_{t \in \T} \sum_{k \in \K} \sum_{(i,j) \in \A}
                c_t \cdot \tij(w_{ikt}) \cdot x_{ijkt}}_{\text{Chi phí thời gian động}}
  \label{eq:obj}
\end{equation}

\noindent
\textbf{Lưu ý về tính phi tuyến (quan trọng):} Thành phần $c_t \cdot \tij(w_{ikt}) \cdot x_{ijkt}$
là tích của biến nhị phân $x_{ijkt}$, biến liên tục $w_{ikt}$, và hàm phi tuyến bậc thang
$\tij(\cdot)$. Do đó, bài toán IRP-TW-DT \textit{về nguyên tắc là MINLP} (Mixed Integer
Nonlinear Program), không phải MILP thuần. Trong nghiên cứu này, tính phi tuyến được xử lý
theo hai cách tùy mục đích:
\begin{enumerate}[noitemsep, label=(\roman*)]
  \item \textit{Trong HGA}: $\tij(\tau)$ được tính trực tiếp bằng thuật toán IGP --- không
        cần linearization.
  \item \textit{Trong MILP lower bound (validation)}: Precompute lookup table $\hat{\tau}_{ij}^{(s)}$
        cho 3 slot xuất phát cố định $\{6\text{h}, 9\text{h}, 19\text{h}\}$, chọn slot tốt
        nhất trước khi giải. Kết quả là MILP tuyến tính có thể giải bằng HiGHS.
\end{enumerate}
Việc tách routing thành \textit{distance cost} và \textit{time-dependent cost} cũng cho
phép phân tích độc lập tác động của giao thông động.

\subsection{Các ràng buộc}

\subsubsection{Ràng buộc tồn kho}

\begin{align}
  I_{it} &= I_{i,t-1} - d_{it} + \sum_{k \in \K} q_{ikt}
             \quad \forall i \in \N',\ t \in \T
             \label{eq:inv_balance} \\
  L_i^{\min} &\leq I_{it} \leq U_i
             \quad \forall i \in \N',\ t \in \T
             \label{eq:inv_bounds} \\
  q_{ikt} &\leq (U_i - I_{i,t-1}) \cdot z_{ikt}
             \quad \forall i \in \N',\ k \in \K,\ t \in \T
             \label{eq:order_up} \\
  q_{ikt} &\geq 0
             \quad \forall i \in \N',\ k \in \K,\ t \in \T
             \label{eq:q_nonneg}
\end{align}

\noindent Ràng buộc (\ref{eq:inv_balance}) là phương trình cân bằng tồn kho.
(\ref{eq:inv_bounds}) đảm bảo tồn kho nằm trong khoảng an toàn.
(\ref{eq:order_up}) thực thi chính sách \textit{Order-up-to-level (OU)}: lượng giao
không vượt dung lượng trống của kho.

\subsubsection{Ràng buộc tải trọng xe}

\begin{equation}
  \sum_{i \in \N'} q_{ikt} \leq Q \quad \forall k \in \K,\ t \in \T
  \label{eq:capacity}
\end{equation}

\subsubsection{Ràng buộc định tuyến (Flow Conservation)}

\begin{align}
  \sum_{j \in \V,\, j \neq i} x_{ijkt} &= z_{ikt}
    \quad \forall i \in \N',\ k \in \K,\ t \in \T
    \label{eq:flow_out} \\
  \sum_{j \in \V,\, j \neq i} x_{jikt} &= z_{ikt}
    \quad \forall i \in \N',\ k \in \K,\ t \in \T
    \label{eq:flow_in} \\
  \sum_{j \in \N'} x_{0jkt} &\leq 1
    \quad \forall k \in \K,\ t \in \T
    \label{eq:depot_out}
\end{align}

\subsubsection{Ràng buộc khung thời gian và giao thông động}

\begin{align}
  w_{jkt} &\geq w_{ikt} + s_i + \tij(w_{ikt}) - M(1 - x_{ijkt})
    \quad \forall (i,j) \in \A,\ k \in \K,\ t \in \T
    \label{eq:time_prop} \\
  e_i \cdot z_{ikt} &\leq w_{ikt}
    \quad \forall i \in \N',\ k \in \K,\ t \in \T
    \label{eq:tw_lb} \\
  w_{ikt} &\leq l_i \cdot z_{ikt} + M(1 - z_{ikt})
    \quad \forall i \in \N',\ k \in \K,\ t \in \T
    \label{eq:tw_ub}
\end{align}

\noindent Ràng buộc (\ref{eq:time_prop}) lan truyền thời gian phục vụ theo travel time
phụ thuộc thời điểm; $w_{0kt}$ là biến tự do đại diện cho giờ xuất phát từ depot.
(\ref{eq:tw_lb})--(\ref{eq:tw_ub}) đảm bảo xe chỉ phục vụ trong khung giờ $[e_i, l_i]$:
khi $z_{ikt}=0$ (không ghé), $w_{ikt}$ không bị ràng buộc; khi $z_{ikt}=1$, buộc
$e_i \leq w_{ikt} \leq l_i$. Lưu ý: $\tij(w_{ikt})$ phi tuyến --- xem xử lý ở mục
\textit{Hàm mục tiêu}.

\subsubsection{Loại bỏ chu trình con (Subtour Elimination)}

\begin{equation}
  \sum_{i \in S} \sum_{j \in S,\, j \neq i} x_{ijkt} \leq |S| - 1
  \quad \forall S \subseteq \N',\ 2 \leq |S| \leq |\N'|,\ k \in \K,\ t \in \T
  \label{eq:sec}
\end{equation}

\noindent Hoặc thay bằng MTZ constraints để giảm số ràng buộc.

\subsubsection{Tính nguyên và không âm}

\begin{equation}
  x_{ijkt} \in \{0,1\},\quad z_{ikt} \in \{0,1\},\quad
  I_{it} \geq 0,\quad q_{ikt} \geq 0,\quad w_{ikt} \geq 0,\quad 0 \leq w_{0kt} \leq H
  \label{eq:integrality}
\end{equation}

% ============================================================
\section{Phương pháp Giải --- Hybrid Genetic Algorithm (HGA)}
% ============================================================

\subsection{Tổng quan Kiến trúc}

IRP-TW-DT là bài toán NP-Hard với không gian tìm kiếm tổ hợp nổ theo hàm mũ.
Thuật toán HGA được thiết kế lại từ cơ chế HGA-CVRPTW của tác giả
(khóa luận trước), bổ sung hai thành phần mới cốt lõi:

\begin{itemize}[leftmargin=*]
  \item \boldterm{Two-part Chromosome}: Mã hóa đồng thời quyết định tồn kho (ma trận
        phân bổ $\mathbf{Y}$) và quyết định định tuyến (chuỗi Giant Tour).
  \item \boldterm{Time-Shift Operator}: Local search trên trục thời gian, dời ngày giao
        của một khách hàng sang ngày có chi phí vận tải thấp hơn mà không vi phạm tồn kho.
\end{itemize}

\subsection{Mã hóa Nghiệm Hai tầng (Two-part Chromosome)}

\begin{equation}
  \text{Chromosome} = \left(\underbrace{\mathbf{Y}_{n \times T}}_{\text{Tầng 1: Allocation}},\ 
                             \underbrace{\boldsymbol{\pi}_{1 \times n}}_{\text{Tầng 2: Giant Tour}}\right)
\end{equation}

\begin{description}[leftmargin=0cm]
  \item[\textbf{Tầng 1 --- Allocation Matrix} $\mathbf{Y} \in \{0,1\}^{n \times T}$:]
    $Y_{it} = 1$ khi khách hàng $i$ được giao hàng vào ngày $t$.
    Lượng hàng tự động được tính theo chính sách OU: $q_{it} = U_i - I_{i,t-1}$.

  \item[\textbf{Tầng 2 --- Giant Tour} $\boldsymbol{\pi}$:]
    Hoán vị của $n$ khách hàng, xác định thứ tự ưu tiên tương đối khi định tuyến.
    Trong mỗi ngày $t$, tập $\{i : Y_{it} = 1\}$ được sắp xếp theo thứ tự của
    $\boldsymbol{\pi}$ rồi đưa vào TD-Split Algorithm.
\end{description}

\subsection{TD-Split Algorithm --- Giải mã và Phân chia Lộ trình}

Với mỗi ngày $t$, chuỗi khách hàng được giao trong ngày đó (ký hiệu $D_t$, sắp xếp
theo Giant Tour) được phân chia thành các route $R_k$ bằng thuật toán quy hoạch động:

\begin{equation}
  V[j] = \min_{i \leq j} \left\{ V[i-1] + C_t(i, j) \right\}
  \quad \forall j = 1, \ldots, |D_t|
\end{equation}

\noindent trong đó $C_t(i,j)$ là chi phí của route phục vụ các khách hàng từ $D_t[i]$
đến $D_t[j]$ trong ngày $t$. Chi phí này được tính theo hàm mục tiêu (\ref{eq:obj})
sử dụng travel time $\tij(\cdot)$ từ mô hình IGP.

\boldterm{Lựa chọn giờ xuất phát:} Để khai thác giao thông động, mỗi ngày, bộ giải thử
3 phương án giờ xuất phát từ depot: $\{6\text{h},\ 9\text{h},\ 19\text{h}\}$ và chọn
phương án có tổng chi phí routing thấp nhất.

\subsection{Toán tử Tiến hóa (Evolutionary Operators)}

\begin{description}[leftmargin=0cm, itemsep=6pt]
  \item[\textbf{Lai ghép (Crossover):}]
    \begin{itemize}[noitemsep]
      \item \textit{Tầng 1}: Uniform Crossover trên ma trận $\mathbf{Y}$.
      \item \textit{Tầng 2}: Order Crossover (OX) trên Giant Tour $\boldsymbol{\pi}$.
    \end{itemize}

  \item[\textbf{Đột biến (Mutation):}]
    \begin{itemize}[noitemsep]
      \item \textit{Tầng 1}: Lật ngẫu nhiên một số bit trong $\mathbf{Y}$, sau đó
            kiểm tra và sửa nếu có khách bị stock-out.
      \item \textit{Tầng 2}: Swap Mutation hoặc Inversion Mutation trên $\boldsymbol{\pi}$.
    \end{itemize}

  \item[\textbf{Repair Operator:}]
    Sau mỗi phép biến đổi, kiểm tra $I_{it} \geq L_i^{\min}$ cho mọi $(i,t)$.
    Nếu vi phạm, buộc $Y_{it} = 1$ và cập nhật lại $q_{it}$.
\end{description}

\subsection{Tìm kiếm Cục bộ (Local Search)}

\begin{description}[leftmargin=0cm, itemsep=6pt]
  \item[\textbf{2-opt (Không gian):}]
    Áp dụng đổi chỗ 2 cung liên tiếp để cải thiện từng route trong mỗi ngày.
    Sử dụng cơ chế đánh giá nhanh của \citealias{visser2020efficient}.

  \item[\textbf{Time-Shift Neighborhood Search (Thời gian) --- Đóng góp chính:}]
    Với mỗi khách hàng $i$ đang được giao vào ngày $t$, thử dời sang ngày $t'$:
    \begin{itemize}[noitemsep]
      \item Tạo $\mathbf{Y}'$: xóa $Y_{it} \leftarrow 0$, đặt $Y_{it'} \leftarrow 1$.
      \item Kiểm tra tính khả thi: $I_{it''} \geq L_i^{\min}$ với mọi $t''$ bị ảnh hưởng.
      \item Nếu khả thi và $Z(\mathbf{Y}') < Z(\mathbf{Y})$: chấp nhận thay đổi.
    \end{itemize}
    Time-Shift neighborhood search cho phép dồn giao hàng sang các ngày có giao thông
    thuận lợi, trực tiếp tối ưu hóa trade-off giữa chi phí lưu kho và chi phí vận tải.
\end{description}

\subsection{Hàm Thích nghi (Fitness Function)}

\begin{equation}
  F(\text{chrom}) = Z(\text{chrom})
                    + \lambda_1 \cdot P_{\text{stock-out}}
                    + \lambda_2 \cdot P_{\text{capacity}}
                    + \lambda_3 \cdot P_{\text{TW}}
\end{equation}

\noindent Trong đó $P_{\cdot}$ là hàm phạt, $\lambda_1 \gg \lambda_2 > \lambda_3 > 0$.
Lời giải với penalty $> 0$ được coi là bất khả thi và được xếp sau lời giải khả thi
bất kể giá trị $Z$.

\newpage
\subsection{Lưu đồ Thuật toán}

\begin{algorithm}[htbp]
\caption{Hybrid Genetic Algorithm cho IRP-TW-DT}
\label{alg:hga}
\begin{algorithmic}[1]
\Require Instance $\mathcal{I}$, $P$ (pop. size), $G$ (generations), $T_{\max}$ (time limit)
\Ensure Chromosome $\text{chrom}^*$ với $F$ tốt nhất

\State Khởi tạo population $\mathcal{P} = \{\text{chrom}_p\}_{p=1}^{P}$ ngẫu nhiên, đảm bảo mỗi khách được giao $\geq 1$ lần
\State Evaluate $F(\text{chrom}_p)$ cho mọi $p$ bằng Decode() + TD-Split
\State $\text{chrom}^* \leftarrow \arg\min_p F(\text{chrom}_p)$

\For{$g = 1$ to $G$}
  \If{thời gian chạy $> T_{\max}$} \textbf{break} \EndIf
  \State $\mathcal{P}' \leftarrow$ Top $10\%$ của $\mathcal{P}$ (elitism)
  \While{$|\mathcal{P}'| < P$}
    \State $p_1, p_2 \leftarrow$ Tournament Selection($\mathcal{P}$, $k=5$)
    \State $c_1, c_2 \leftarrow$ Crossover$(p_1, p_2)$ \Comment{OX + Uniform Crossover}
    \State Mutate$(c_1)$; Mutate$(c_2)$
    \State Repair$(c_1, \mathcal{I})$; Repair$(c_2, \mathcal{I})$
    \State Evaluate $F(c_1)$, $F(c_2)$
    \State $\mathcal{P}' \leftarrow \mathcal{P}' \cup \{c_1, c_2\}$
  \EndWhile
  \State $\mathcal{P} \leftarrow$ Sort$(\mathcal{P}')$, giữ $P$ cá thể tốt nhất
  \For{$p = 1$ to $\lceil P/5 \rceil$} \Comment{Local search top 20\%}
    \State $\mathcal{P}[p] \leftarrow$ 2opt$(\mathcal{P}[p])$
    \State $\mathcal{P}[p] \leftarrow$ TimeShiftNS$(\mathcal{P}[p], \mathcal{I})$
  \EndFor
  \If{$F(\mathcal{P}[1]) < F(\text{chrom}^*)$}
    $\text{chrom}^* \leftarrow \mathcal{P}[1]$
  \EndIf
\EndFor
\State \Return $\text{chrom}^*$
\end{algorithmic}
\end{algorithm}

% ============================================================
\section{Dữ liệu Nghiên cứu}
% ============================================================

\subsection{Bộ dữ liệu mô phỏng Hà Nội (Synthetic Hanoi Dataset)}

Dữ liệu được sinh tổng hợp, kế thừa cấu trúc phân cụm không gian từ khóa luận CVRPTW
với $K=5$ cụm phân bố trong bán kính 15 km từ depot (trung tâm Hoàn Kiếm). Các tham số
tồn kho được sinh ngẫu nhiên:

\begin{itemize}[leftmargin=*]
  \item $U_i \sim \mathcal{U}(80, 200)$ đơn vị.
  \item $L_i^{\min} = U_i \cdot \mathcal{U}(0.10, 0.20)$ (tồn kho an toàn 10--20\%).
  \item $I_{i0} \sim \mathcal{U}(L_i^{\min}, 0.6 \cdot U_i)$ (bắt đầu ở mức 30--60\%).
  \item $d_{it} \sim \text{Lognormal}(\mu_d, 0.3)$ với $\mu_d = \log(U_i/7)$
        (nhu cầu trung bình $\approx U_i/7 \approx 14\%$ dung lượng/ngày ---
        phù hợp văn học IRP; đủ dư địa để tối ưu hóa lịch giao trong 7 ngày).
  \item $h_i \sim \mathcal{U}(100, 500)$ VND/đơn vị/ngày.
  \item Khung thời gian: 2 ca (8h--12h hoặc 14h--18h), gán ngẫu nhiên cho từng khách.
  \item $s_i = 0{,}25$h (hằng số cho mọi khách hàng).
\end{itemize}

\noindent\textbf{Giả định tính khả thi (Single-visit Feasibility):}
Mỗi instance được sinh đảm bảo điều kiện: mỗi khách hàng $i$ có thể được phục vụ đơn
lẻ trong khung thời gian của mình nếu xe xuất phát từ depot vào đầu ca (6h hoặc 13h).
Cụ thể, với tốc độ tệ nhất $v^{(2)} = 15$~km/h và bán kính tối đa 15~km, thời gian đến
khách xa nhất là $15/15 = 1$~giờ, đủ để vào cửa sổ [8h, 12h] nếu xuất phát lúc 6h.
Điều kiện này được kiểm tra tự động trong quá trình sinh dữ liệu:
$e_i \geq w_{0k} + \tau_{0i}(w_{0k}) + s_i$ với ít nhất một giờ xuất phát khả thi.
Nếu vi phạm, instance được sinh lại.

\subsection{Tham số hệ thống}

\begin{center}
\begin{tabular}{lll}
\toprule
Tham số & Giá trị & Nguồn \\
\midrule
Chân trời kế hoạch $T$ & 7 ngày & Tiêu chuẩn IRP văn học \\
Số xe $m$ & 2/3/5 (tương ứng S/M/L) & Đề xuất \\
Tải trọng xe $Q$ & 500 đơn vị & Đề xuất \\
$c_d$ (chi phí khoảng cách) & 3.500 VND/km & Ước tính: chi phí nhiên liệu + bảo trì biến đổi \\
                              &               & xe tải nhỏ; \textit{cần hiệu chỉnh thực địa} \\
$c_t$ (chi phí thời gian) & 68.000--80.000 VND/giờ & Từ mức lương tài xế Hà Nội 2024 \\
                           &                          & 8--15 triệu VND/tháng \citep{jobsgo2024} \\
Speed profile & Bảng \ref{tab:speed} & TomTom Traffic Index 2025 \\
\bottomrule
\end{tabular}
\end{center}

\subsection{Bộ Benchmark Chuẩn (để Validate HGA)}

Trước khi chạy trên dữ liệu Hà Nội, HGA được kiểm chứng trên các instance chuẩn của
\citealias{archetti2007branch} (bộ Archetti instances, 5--50 khách) để so sánh với
kết quả exact solver. Bộ benchmark mới của \citealias{skalnes2024benchmark}
(270 instances, 10--200 khách) cũng sẽ được xem xét để kiểm tra scalability.

% ============================================================
\section{Thiết kế Thực nghiệm}
% ============================================================

\subsection{Bốn Kịch bản So sánh}

\begin{table}[htbp]
\centering
\caption{Thiết kế bốn kịch bản thực nghiệm}
\label{tab:scenarios}
\renewcommand{\arraystretch}{1.3}
\begin{tabularx}{\linewidth}{@{}l l X@{}}
\toprule
\textbf{Kịch bản} & \textbf{Mô hình} & \textbf{Mô tả} \\
\midrule
\rowcolor{lightblue}
P --- Periodic & Heuristic & Giao hàng định kỳ mỗi 3 ngày cho tất cả khách.
           Phản ánh cách vận hành đơn giản nhất trong thực tế.
           Tốc độ hằng $\bar{v} = 18$~km/h. \\
\rowcolor{lightred}
A --- RMI & CVRPTW & Xe giao ngay ngày hôm sau khi khách báo
           $I_{it} \leq L_i^{\min}$. Tốc độ hằng $\bar{v} = 18$~km/h.
           Phản ánh thực tế phổ biến của doanh nghiệp logistics Việt Nam. \\
\rowcolor{lightyellow}
B --- IRP-TW-Static & IRP-TW & Mô hình IRP-TW nhưng tắt giao thông động,
           dùng cùng tốc độ hằng $\bar{v} = 18$~km/h.
           So sánh B vs.\ A: đóng góp thuần của VMI scheduling;
           so sánh B vs.\ C: đóng góp thuần của Dynamic Traffic. \\
\rowcolor{lightgreen}
C --- IRP-TW-DT & IRP-TW-DT & Mô hình đề xuất đầy đủ với profile giao thông
           Hà Nội 5 múi giờ và HGA kết hợp Time-Shift neighborhood search. \\
\bottomrule
\end{tabularx}
\end{table}

\subsection{Quy mô Thực nghiệm}

\begin{itemize}[leftmargin=*]
  \item \textbf{Kích thước}: 4 kịch bản $\times$ 3 quy mô $\times$ 5 seeds = 60 runs.
  \item \textbf{Quy mô nhỏ (S)}: $n = 20$, $m = 2$ xe.
  \item \textbf{Quy mô trung bình (M)}: $n = 50$, $m = 3$ xe.
  \item \textbf{Quy mô lớn (L)}: $n = 100$, $m = 5$ xe.
  \item \textbf{Seeds}: $\{42, 123, 456, 789, 1000\}$ để đảm bảo tính tái lập.
  \item \textbf{Thời gian giới hạn HGA}: 5 phút/instance với NumPy vectorization
        (ước tính cho $n=100$: $\sim$1--2 phút; Python thuần: $\sim$10--30 phút).
  \item \textbf{Kịch bản P (Periodic)}: không dùng HGA; giao cố định mỗi 3 ngày,
        routing giải bằng nearest-neighbor heuristic --- implement trong $< 1$ giờ.
\end{itemize}

\subsection{Chỉ số Đánh giá (KPIs)}

\begin{table}[htbp]
\centering
\caption{Các chỉ số đánh giá hiệu năng}
\label{tab:kpis}
\renewcommand{\arraystretch}{1.25}
\begin{tabularx}{\linewidth}{@{}l l X@{}}
\toprule
\textbf{KPI} & \textbf{Đơn vị} & \textbf{Ý nghĩa và cách tính} \\
\midrule
Total Cost ($Z$) & VND & Hàm mục tiêu (\ref{eq:obj}): tổng lưu kho + routing \\
Inventory Cost   & VND & $\sum_t \sum_i h_i I_{it}$ --- đo mức ``găm hàng'' \\
Routing Cost     & VND & Tổng chi phí khoảng cách và thời gian \\
Avg. Inventory Level & \% $U_i$ & Mức tồn kho trung bình, phản ánh trade-off \\
Total Distance   & km & Tổng quãng đường tất cả xe \\
TW Compliance Rate & \% & Tỷ lệ chuyến giao đúng trong time window \\
\# Deliveries    & --  & Tổng số lần giao trên toàn bộ horizon \\
CPU Time         & giây & Thời gian tính toán \\
\bottomrule
\end{tabularx}
\end{table}

\subsection{Phương pháp Phân tích Kết quả}

\begin{itemize}[leftmargin=*]
  \item \boldterm{P vs. A (Periodic vs. RMI)}: Kiểm tra xem reactive demand-driven
        có thực sự tốt hơn giao định kỳ không --- kết quả có thể bất ngờ.
  \item \boldterm{A vs. B (RMI vs. IRP-TW-Static)}: Tính phần trăm tiết kiệm khi
        chuyển từ reactive sang VMI scheduling với tốc độ hằng. Kiểm định t-test.
  \item \boldterm{B vs. C (Static vs. Dynamic)}: Phân tích đóng góp thuần của
        dynamic traffic bằng kiểm định Wilcoxon signed-rank.
  \item \boldterm{Phân tích Trade-off}: Scatter plot (Inventory Cost, Routing Cost)
        phân tầng theo kịch bản --- kiểm tra giả thuyết \textit{"giao thông tắc nghẽn
        buộc nhà cung cấp chấp nhận tồn kho cao hơn để tránh giờ cao điểm"}.
  \item \boldterm{Validation HGA}: So sánh gap \% giữa HGA và MILP lower bound (HiGHS)
        trên instances nhỏ ($n \leq 10$, $T = 3$, tốc độ cố định). Mục tiêu: gap $\leq 15\%$.
        Lưu ý: với $n=20$ và $T=7$, MILP có thể không hội tụ trong 10 phút ---
        validate chỉ thực hiện trên tập con nhỏ.
  \item \boldterm{Convergence analysis}: Vẽ đường hội tụ fitness theo số generations.
\end{itemize}

% ============================================================
\section{Đóng góp Dự kiến}
% ============================================================

\begin{description}[leftmargin=0cm, itemsep=8pt]
  \item[\boldterm{Đóng góp lý luận:}]
    Mô hình IRP-TW-DT là mở rộng của TD-IRP \citep{touzout2022assign} với ràng buộc
    time windows --- khoảng trống chưa được lấp. Công trình cung cấp bằng chứng
    định lượng lần đầu tiên về trade-off giữa \textit{chi phí lưu kho} và
    \textit{chi phí ách tắc đô thị} khi tích hợp time windows vào IRP.

  \item[\boldterm{Đóng góp phương pháp:}]
    Thuật toán HGA với Two-part Chromosome là một trong số ít metaheuristic
    giải TD-IRP trong văn học (sau matheuristic của Touzout et al., 2022).
    \textit{Time-Shift neighborhood search} --- tìm kiếm lân cận trên trục thời gian ---
    là cơ chế mới khác về bản chất với 2-opt tối ưu định tuyến trên không gian.

  \item[\boldterm{Đóng góp thực tiễn:}]
    Kết quả thực nghiệm cung cấp cơ sở ra quyết định (Decision Support) cho doanh
    nghiệp logistics tại Hà Nội khi cân nhắc chuyển sang VMI. Cụ thể, nghiên cứu
    kiểm tra giả thuyết: \textit{tắc nghẽn đô thị có thể justify việc giao hàng sớm
    hơn và duy trì tồn kho cao hơn} --- trade-off mà mô hình tĩnh không thể định
    lượng. Insight này trực tiếp ảnh hưởng đến quyết định về buffer stock và
    lịch giao hàng trong thực tế.

  \item[\boldterm{Đóng góp dữ liệu:}]
    Công bố bộ synthetic benchmark Hà Nội (mã nguồn mở) để khuyến khích nghiên cứu
    IRP trong bối cảnh đô thị đang phát triển của châu Á.
\end{description}

% ============================================================
\section{Kế hoạch Thực hiện}
% ============================================================

\begin{table}[htbp]
\centering
\caption{Roadmap thực hiện (10 tuần)}
\label{tab:roadmap}
\renewcommand{\arraystretch}{1.3}
\begin{tabularx}{\linewidth}{@{}c X X@{}}
\toprule
\textbf{Tuần} & \textbf{Công việc} & \textbf{Kết quả kiểm chứng} \\
\midrule
1 & Đọc Touzout et al. (2022) toàn văn; cài môi trường Python + NumPy;
    implement IGP traffic model. &
    Unit test FIFO: $\forall \tau_1 < \tau_2$: arrival$(\tau_1) \leq$ arrival$(\tau_2)$. \\
2 & Thiết kế data structures (dataclasses); implement inventory simulation;
    sinh bộ dữ liệu Hà Nội cho $n = 20$. &
    Kiểm tra cân bằng tồn kho, phát hiện bug tồn kho âm. \\
3 & TD-Split Algorithm; tính $C_t(i,j)$ với time-dependent cost;
    precompute lookup table $\hat{\tau}_{ij}^{(s)}$ cho MILP. &
    Split không vi phạm tải trọng; cost tính đúng trên test case nhỏ. \\
4 & Full fitness evaluation; Two-part Chromosome; random initialization;
    chạy single-thread GA trên $n=20$. &
    Evaluate() chạy đúng; penalty hoạt động; fitness giảm nhẹ sau khởi tạo. \\
5 & Crossover (OX + Uniform), Mutation, Repair operator hoàn chỉnh. &
    Fitness giảm rõ qua generations trên $n=20$, 5 seeds. \\
6 & Time-Shift operator; 2-opt với fast evaluation;
    tuning siêu tham số (pop, mutation rate). &
    Time-Shift cải thiện $\geq 3\%$ so với không có local search ($n=20$). \\
7 & Cài HiGHS Python API; code MILP lower bound ($n \leq 10$, $T=3$);
    validate HGA trên tập nhỏ này. &
    MILP hội tụ $< 5$ phút với $n=10$; HGA gap $\leq 15\%$ là mục tiêu. \\
8 & Implement Periodic Delivery baseline (nearest-neighbor + fixed schedule);
    sinh dữ liệu cho $n = 50, 100$; chạy 60 experiments
    (4 kịch bản $\times$ 3 quy mô $\times$ 5 seeds). &
    Tất cả 60 runs hoàn thành; kết quả lưu CSV. \\
9 & Phân tích thống kê (t-test, Wilcoxon); scatter plot trade-off;
    convergence curves. &
    Kiểm định có $p < 0.05$ cho ít nhất 1 so sánh có ý nghĩa. \\
10 & Viết báo cáo hoàn chỉnh; review nội bộ; chuẩn bị slide. &
    Bản thảo báo cáo với đủ bảng kết quả, hình, và thảo luận. \\
\bottomrule
\end{tabularx}
\end{table}

% ============================================================
\section{Kết luận}
% ============================================================

Nghiên cứu đề xuất mô hình IRP-TW-DT --- bài toán Inventory Routing Problem tích hợp
Time Windows và Giao thông Động --- như một bước tiến tự nhiên tiếp theo sau công trình
TD-IRP của \citealias{touzout2022assign}. Bốn kịch bản thực nghiệm (Periodic, RMI,
IRP-TW-Static, IRP-TW-DT) cho phép tách biệt đóng góp của từng thành phần: VMI scheduling,
time windows, và dynamic traffic. Time-Shift neighborhood search trên trục thời gian là
cơ chế thuật toán phân biệt rõ đề tài này với các công trình trước. Kết quả dự kiến sẽ
định lượng được mức độ tắc nghẽn đô thị ảnh hưởng đến trade-off tồn kho -- vận tải,
cung cấp nền tảng khoa học thực tiễn cho logistics chặng cuối tại các đô thị đang phát
triển nhanh của châu Á.

% ============================================================
%  REFERENCES
% ============================================================
\newpage
\bibliographystyle{abbrvnat}

\begin{thebibliography}{99}

\bibitem[Adamo et al.(2024)]{adamo2024review}
Adamo, T., Gendreau, M., Ghiani, G., \& Guerriero, E. (2024).
A review of recent advances in time-dependent vehicle routing.
\textit{European Journal of Operational Research}, 319(1), 1--15.

\bibitem[Adulyasak et al.(2014)]{adulyasak2014branch}
Adulyasak, Y., Cordeau, J.-F., \& Jans, R. (2014).
Formulations and branch-and-cut algorithms for multivehicle production and inventory routing problems.
\textit{INFORMS Journal on Computing}, 26(1), 103--120.

\bibitem[Archetti et al.(2007)]{archetti2007branch}
Archetti, C., Bertazzi, L., Laporte, G., \& Speranza, M. G. (2007).
A branch-and-cut algorithm for a vendor-managed inventory-routing problem.
\textit{Transportation Science}, 41(3), 382--391.

\bibitem[Archetti \& Speranza(2016)]{archetti2016vmi}
Archetti, C., \& Speranza, M. G. (2016).
The inventory routing problem: The value of integration.
\textit{International Transactions in Operational Research}, 23(3), 393--407.

\bibitem[Bouvier et al.(2024)]{bouvier2024continent}
Bouvier, T., Vidal, T., Gendreau, M., \& Laporte, G. (2024).
Solving a continent-scale inventory routing problem at Renault.
\textit{Transportation Science}, 58(2), 340--360.

\bibitem[Coelho et al.(2014)]{coelho2014thirty}
Coelho, L. C., Cordeau, J.-F., \& Laporte, G. (2014).
Thirty years of inventory routing.
\textit{Transportation Science}, 48(1), 1--19.

\bibitem[Desaulniers et al.(2016)]{desaulniers2016branch}
Desaulniers, G., Rakke, J. G., \& Coelho, L. C. (2016).
A branch-price-and-cut algorithm for the inventory routing problem.
\textit{Transportation Science}, 50(3), 1060--1076.

\bibitem[Gevaers et al.(2011)]{gevaers2011characterization}
Gevaers, R., Van de Voorde, E., \& Vanelslander, T. (2011).
Characteristics and typology of last-mile logistics from an innovation perspective in an urban context.
\textit{City Distribution and Urban Freight Transport: Multiple Perspectives}, 56--71.

\bibitem[FedEx(2024)]{fedex2024apac}
FedEx Vietnam. (2024).
Best Last-Mile Delivery Options for APAC SMEs.
Retrieved from \url{https://www.fedex.com/en-vn/business-insights/ecommerce/last-mile-delivery-apac-smes.html}

\bibitem[Allied Market Research(2022)]{amr2022vietnam}
Allied Market Research. (2022).
Vietnam Express Delivery Services Market.
Retrieved from \url{https://www.alliedmarketresearch.com/vietnam-express-delivery-services-market-A11094}

\bibitem[TomTom(2025)]{tomtom2025hanoi}
TomTom. (2025).
Ha Noi Traffic Report --- TomTom Traffic Index 2025.
Retrieved from \url{https://www.tomtom.com/traffic-index/city/ha-noi}

\bibitem[JobsGO(2024)]{jobsgo2024}
JobsGO. (2024).
Mức lương tài xế hiện nay.
Retrieved from \url{https://jobsgo.vn/muc-luong-tai-xe.html}

\bibitem[Ghiani \& Guerriero(2014)]{ghiani2014note}
Ghiani, G., \& Guerriero, E. (2014).
A note on the Ichoua, Gendreau, and Potvin (2003) travel time model.
\textit{Transportation Science}, 48(3), 458--462.

\bibitem[Ichoua et al.(2003)]{ichoua2003vehicle}
Ichoua, S., Gendreau, M., \& Potvin, J.-Y. (2003).
Vehicle dispatching with time-dependent travel times.
\textit{European Journal of Operational Research}, 144(2), 379--396.

\bibitem[Rifki et al.(2020)]{rifki2020impact}
Rifki, O., Chiabaut, N., \& Solnon, C. (2020).
On the impact of spatio-temporal granularity of traffic conditions on the quality of pickup and delivery optimal tours.
\textit{Transportation Research Part E}, 142, 102043.

\bibitem[Skålnes et al.(2022)]{skalnes2022improved}
Skålnes, J., Andersson, H., Desaulniers, G., \& Stålhane, M. (2022).
An improved formulation for the inventory routing problem with time-varying demands.
\textit{European Journal of Operational Research}, 302(3), 1189--1201.

\bibitem[Skålnes et al.(2024)]{skalnes2024benchmark}
Skålnes, J., Ben Ahmed, M., Hvattum, L. M., \& Stålhane, M. (2024).
New benchmark instances for the inventory routing problem.
\textit{European Journal of Operational Research}, 313(3), 959--975.

\bibitem[Touzout et al.(2022)]{touzout2022assign}
Touzout, F. A., Ladier, A.-L., \& Hadj-Hamou, K. (2022).
An assign-and-route matheuristic for the time-dependent inventory routing problem.
\textit{European Journal of Operational Research}, 300(3), 1081--1097.

\bibitem[Vidal et al.(2012)]{vidal2012hga}
Vidal, T., Crainic, T. G., Gendreau, M., \& Prins, C. (2012).
A hybrid genetic algorithm for multidepot and periodic vehicle routing problems.
\textit{Operations Research}, 60(3), 611--624.

\bibitem[Visser \& Spliet(2020)]{visser2020efficient}
Visser, T. R., \& Spliet, R. (2020).
Efficient move evaluations for time-dependent vehicle routing problems.
\textit{Transportation Science}, 54(4), 1091--1112.

\end{thebibliography}

\end{document}